% Methodological Appendix for Tacet USENIX Security 2026 Paper
% To be included via % Methodological Appendix for Tacet USENIX Security 2026 Paper
% To be included via % Methodological Appendix for Tacet USENIX Security 2026 Paper
% To be included via % Methodological Appendix for Tacet USENIX Security 2026 Paper
% To be included via \input{appendix_methodology} in the main document

\section{Statistical Methodology Details}
\label{app:methodology}

This appendix provides additional detail on tacet's statistical methodology for reviewers familiar with Bayesian inference. We cover the probabilistic model, inference procedure, autocorrelation handling, quality gates, and measurement floor derivation. For complete algorithmic specifications, see the specification document in the artifact repository.

\subsection{Probabilistic Model}
\label{app:model}

\paragraph{Test Statistic.}
Rather than comparing means (as in t-tests), we compare timing distributions via their deciles. Let $F$ and $R$ denote timing samples from fixed (e.g., all-zeros) and random input classes, respectively. The test statistic is the vector of quantile differences:
\begin{equation}
\Delta = \hat{q}(F) - \hat{q}(R) \in \mathbb{R}^9
\end{equation}
where $\hat{q}_p$ denotes the empirical $p$-th quantile for $p \in \{0.1, 0.2, \ldots, 0.9\}$. This captures both uniform shifts (all quantiles affected equally, indicating branch changes) and tail effects (upper quantiles affected more, indicating cache-based leaks).

The decision functional is $m(\delta) = \max_k |\delta_k|$---the maximum absolute quantile difference. A timing leak exceeding threshold $\theta$ is declared when $m(\delta) > \theta$.

\paragraph{Likelihood.}
Quantile estimators are asymptotically normal, but real timing data can deviate from this approximation when covariance is underestimated. To prevent pathological certainty under covariance misspecification, we use a robust Student-$t$ likelihood via a scale-mixture representation:
\begin{align}
\Delta \mid \delta, \kappa &\sim \mathcal{N}\!\left(\delta, \frac{\Sigma_n}{\kappa}\right) \\
\kappa &\sim \text{Gamma}\!\left(\frac{\nu_\ell}{2}, \frac{\nu_\ell}{2}\right)
\end{align}
with $\nu_\ell = 8$. Marginally, this yields $\Delta \mid \delta \sim t_{\nu_\ell}(\delta, \Sigma_n)$.

When the estimated covariance $\Sigma_n$ is too small (e.g., due to unmodeled autocorrelation), $\kappa$ is pulled below 1, automatically inflating uncertainty. This is more robust than a Gaussian likelihood, which would produce overconfident posteriors under covariance misestimation.

\paragraph{Prior.}
We place a multivariate Student-$t$ prior on the latent effect profile $\delta$:
\begin{equation}
\delta \sim t_\nu(0, \sigma^2 R), \quad \nu = 4
\end{equation}
where $R = \text{Corr}(\Sigma_{\text{rate}})$ is the prior correlation matrix derived from the covariance rate estimated during calibration. The scale $\sigma$ is calibrated so that the prior probability of exceeding the effective threshold $\theta_{\text{eff}}$ equals a target $\pi_0 = 0.62$:
\begin{equation}
P\!\left(\max_k |\delta_k| > \theta_{\text{eff}} \mid \delta \sim t_\nu(0, \sigma^2 R)\right) = \pi_0
\end{equation}
This is solved by deterministic 1D root-finding using Monte Carlo integration (50,000 samples). The resulting prior expresses genuine uncertainty: it does not strongly favor either ``leak'' or ``no leak'' before seeing data.

We implement inference using the equivalent scale-mixture representation:
\begin{align}
\lambda &\sim \text{Gamma}\!\left(\frac{\nu}{2}, \frac{\nu}{2}\right) \\
\delta \mid \lambda &\sim \mathcal{N}\!\left(0, \frac{\sigma^2}{\lambda} R\right)
\end{align}
which enables closed-form Gibbs conditionals.

\subsection{Inference Procedure}
\label{app:inference}

\paragraph{Gibbs Sampler.}
We approximate the posterior $p(\delta \mid \Delta)$ using a deterministic Gibbs sampler over $(\delta, \lambda, \kappa)$. Each iteration samples:
\begin{enumerate}
\item $\delta^{(t)} \sim p(\delta \mid \lambda^{(t-1)}, \kappa^{(t-1)}, \Delta)$ --- a multivariate normal
\item $\lambda^{(t)} \sim p(\lambda \mid \delta^{(t)})$ --- a gamma distribution
\item $\kappa^{(t)} \sim p(\kappa \mid \delta^{(t)}, \Delta)$ --- a gamma distribution
\end{enumerate}

All conditionals are conjugate and admit closed-form sampling. The $\delta$ conditional requires computing a posterior precision matrix $Q = \kappa \Sigma_n^{-1} + (\lambda/\sigma^2) R^{-1}$ and sampling from $\mathcal{N}(Q^{-1}(\kappa \Sigma_n^{-1} \Delta), Q^{-1})$. All matrix operations use Cholesky decomposition for numerical stability---we never form explicit inverses.

\paragraph{Convergence and Cost.}
We run 256 total iterations with 64 burn-in, retaining 192 samples. This is sufficient because: (1) the model is low-dimensional ($d=9$) with conjugate structure, ensuring fast mixing; (2) initialization at $\lambda^{(0)} = \kappa^{(0)} = 1$ (the prior means) provides a good starting point; and (3) the decision (leak probability) depends only on whether $m(\delta^{(s)}) > \theta_{\text{eff}}$, which is robust to modest posterior approximation error.

We track $\lambda$ mixing via its coefficient of variation (CV) and effective sample size (ESS). If $\lambda$ CV $< 0.1$ or $\lambda$ ESS $< 20$, we flag a potential convergence issue in diagnostics, though this rarely occurs in practice.

Computational cost per posterior update is $O(d^3 + d^2 \cdot N_{\text{gibbs}})$, dominated by the Cholesky decompositions. With $d=9$ and $N_{\text{gibbs}}=256$, this takes $<1$ms on modern hardware---negligible compared to measurement time.

\paragraph{Leak Probability.}
The leak probability is estimated as the Monte Carlo fraction of posterior samples exceeding the threshold:
\begin{equation}
\widehat{P}(\text{leak} > \theta_{\text{eff}} \mid \Delta) = \frac{1}{N_{\text{keep}}} \sum_{s=1}^{N_{\text{keep}}} \mathbf{1}\!\left[m(\delta^{(s)}) > \theta_{\text{eff}}\right]
\end{equation}

\subsection{Block Bootstrap for Autocorrelation}
\label{app:bootstrap}

\paragraph{Why Standard Bootstrap Fails.}
Timing measurements exhibit autocorrelation: nearby samples are more similar than distant ones due to cache state, CPU frequency scaling, and OS scheduling. Standard bootstrap assumes i.i.d.\ samples, systematically underestimating variance and producing overconfident posteriors.

For example, if consecutive measurements have autocorrelation $\rho = 0.3$, the effective sample size may be 30--50\% lower than the nominal count. Ignoring this inflates the false positive rate---the main failure mode we aim to prevent.

\paragraph{Stream-Based Block Bootstrap.}
We use block bootstrap on the \emph{acquisition stream}---the interleaved sequence of (class, timing) pairs in measurement order---rather than per-class samples. This preserves the true temporal dependence structure.

The algorithm resamples contiguous blocks of length $\hat{b}$ with replacement, reconstructs a bootstrap stream of the same length, splits by class, and computes quantile differences. Over $B = 2000$ iterations, we estimate the covariance matrix $\Sigma_{\text{cal}}$ of $\Delta$.

\paragraph{Block Length Selection.}
We use the Politis-White automatic block length selection algorithm with a class-conditional modification. Standard Politis-White estimates autocorrelation from raw lag-$k$ pairs, but with interleaved sampling, a lag-$k$ pair may cross class boundaries. Instead, we compute autocorrelation using only same-class pairs at each acquisition-stream lag, then take the conservative maximum across classes.

The optimal block length scales as $\hat{b} \propto T^{1/3}$ where $T$ is the stream length. We clamp $\hat{b} \in [10, \min(3\sqrt{T}, T/3)]$ to prevent degenerate blocks. In high-autocorrelation regimes (detected when $\rho_{\text{max}}(10) > 0.3$), we inflate the block length by 50\% for additional conservatism.

\paragraph{Effective Sample Size via IACT.}
We compute the integrated autocorrelation time (IACT) $\hat{\tau}$ using the Geyer initial monotone sequence (IMS) algorithm. This estimates how many nominal samples equal one independent sample:
\begin{equation}
n_{\text{eff}} = \frac{n}{\hat{\tau}}
\end{equation}

For timing analysis, IACT is computed on indicator series $\mathbf{1}[y \leq q_p]$ for each quantile $p$, taking the maximum across quantiles and classes. This captures dependence in the quantile estimators, not just raw timings.

\paragraph{Covariance Scaling.}
The key insight enabling adaptive sampling is that quantile estimator covariance scales as $1/n$. We estimate a ``covariance rate'' during calibration:
\begin{equation}
\Sigma_{\text{rate}} = \hat{\Sigma}_{\text{cal}} \cdot n_{\text{cal}}
\end{equation}
During the adaptive loop with $n$ samples and current IACT $\hat{\tau}$:
\begin{equation}
\Sigma_n = \Sigma_{\text{rate}} \cdot \frac{\hat{\tau}}{\hat{\tau}_{\text{cal}}} \cdot \frac{1}{n}
\end{equation}
This allows cheap posterior updates without re-bootstrapping at each batch.

\subsection{Quality Gates}
\label{app:gates}

Quality gates detect when data quality is insufficient for a confident verdict, forcing an \textsc{Inconclusive} outcome rather than an unreliable Pass or Fail. This is the key mechanism ensuring calibrated verdicts.

\paragraph{Gate 1: Insufficient Information Gain.}
We compute the KL divergence between Gaussian approximations of the prior and posterior:
\begin{multline}
\text{KL}(p_{\text{post}} \| p_{\text{prior}}) = \frac{1}{2}\Big(
\text{tr}(\Lambda_0^{-1}\Lambda_{\text{post}}) \\
+ \mu_{\text{post}}^\top \Lambda_0^{-1} \mu_{\text{post}}
- d + \ln\frac{|\Lambda_0|}{|\Lambda_{\text{post}}|}
\Big)
\end{multline}
where $\Lambda_0 = 2\sigma^2 R$ is the marginal prior covariance for $\nu=4$, and $(\mu_{\text{post}}, \Lambda_{\text{post}})$ are estimated from Gibbs samples.

\textbf{DataTooNoisy}: If KL $< 0.7$ nats after calibration, the posterior has barely moved from the prior---the data is too noisy to distinguish signal from noise.

\textbf{NotLearning}: If the sum of inter-batch KL divergences falls below 0.001 for 5+ consecutive batches, the posterior has stopped updating. Additional samples will not help.

\paragraph{Gate 2: Would Take Too Long.}
We extrapolate time to decision based on current convergence rate. If projected time exceeds 10$\times$ the configured budget, we stop early with \textsc{Inconclusive} rather than waste resources on a test that cannot conclude.

\paragraph{Gate 3: Resource Budget Exceeded.}
If elapsed time exceeds the configured time budget, or total samples exceed the configured maximum, we stop with \textsc{Inconclusive} and report the current posterior probability.

\paragraph{Gate 4: Condition Drift Detected.}
The covariance rate $\Sigma_{\text{rate}}$ is estimated during calibration. If measurement conditions change during the adaptive loop, this estimate becomes invalid. We detect drift by comparing:
\begin{itemize}
\item Variance ratio: $\sigma_{\text{post}}^2 / \sigma_{\text{cal}}^2$ outside $[0.5, 2.0]$
\item Autocorrelation change: $|\rho_{\text{post}}(1) - \rho_{\text{cal}}(1)| > 0.3$
\item Mean drift: $|\mu_{\text{post}} - \mu_{\text{cal}}| / \sigma_{\text{cal}} > 3.0$
\end{itemize}
Any trigger yields \textsc{Inconclusive} with reason \textsc{ConditionsChanged}.

\paragraph{Likelihood Inflation Warning.}
When the posterior mean of $\kappa$ falls below 0.3, the likelihood covariance was effectively scaled up by $>3\times$ for robustness. This indicates potential covariance misestimation and is reported in diagnostics, though it does not block verdicts.

\subsection{Measurement Floor Derivation}
\label{app:floor}

A critical design element is distinguishing what the user \emph{wants} to detect ($\theta_{\text{user}}$) from what the measurement \emph{can} detect ($\theta_{\text{floor}}$).

\paragraph{Floor-Rate Constant.}
During calibration, we compute a floor-rate constant $c_{\text{floor}}$ representing the 95th percentile of the null distribution of $m(\delta)$:
\begin{equation}
c_{\text{floor}} = q_{0.95}\!\left(\max_k |Z_{0,k}|\right), \quad Z_0 \sim \mathcal{N}(0, \Sigma_{\text{rate}})
\end{equation}
estimated via Monte Carlo (50,000 samples). This characterizes the ``noise floor'' of the measurement system.

\paragraph{Dynamic Measurement Floor.}
During the adaptive loop, the statistical floor shrinks with effective sample size:
\begin{equation}
\theta_{\text{floor,stat}}(n) = \frac{c_{\text{floor}}}{\sqrt{n_{\text{eff}}}}
\end{equation}

The timer also imposes a floor based on quantization:
\begin{equation}
\theta_{\text{tick}} = \frac{\text{1 tick (ns)}}{K}
\end{equation}
where $K$ is the batch size (if batching is used). The combined floor is:
\begin{equation}
\theta_{\text{floor}}(n) = \max\!\left(\theta_{\text{floor,stat}}(n), \theta_{\text{tick}}\right)
\end{equation}

\paragraph{Effective Threshold.}
The effective threshold used for inference is:
\begin{equation}
\theta_{\text{eff}} = \max(\theta_{\text{user}}, \theta_{\text{floor}})
\end{equation}

When $\theta_{\text{eff}} > \theta_{\text{user}}$:
\begin{itemize}
\item \textbf{Fail propagates}: Detecting $m(\delta) > \theta_{\text{eff}}$ implies $m(\delta) > \theta_{\text{user}}$, since $\theta_{\text{eff}} \geq \theta_{\text{user}}$.
\item \textbf{Pass does not propagate}: ``No effect above $\theta_{\text{eff}}$'' is compatible with effects in $(\theta_{\text{user}}, \theta_{\text{eff}}]$. We return \textsc{Inconclusive} with reason \textsc{ThresholdElevated}.
\end{itemize}

As sampling continues, $\theta_{\text{floor}}(n)$ decreases. If it drops to $\theta_{\text{user}}$ or below, Pass becomes possible (subject to the posterior). This is why adaptive sampling can eventually reach Pass even when early calibration shows an elevated floor.

\vspace{1em}
\noindent\textit{For complete algorithmic details including pseudocode, numerical stability procedures, and platform-specific considerations, see the specification document in the artifact repository.}
 in the main document

\section{Statistical Methodology Details}
\label{app:methodology}

This appendix provides additional detail on tacet's statistical methodology for reviewers familiar with Bayesian inference. We cover the probabilistic model, inference procedure, autocorrelation handling, quality gates, and measurement floor derivation. For complete algorithmic specifications, see the specification document in the artifact repository.

\subsection{Probabilistic Model}
\label{app:model}

\paragraph{Test Statistic.}
Rather than comparing means (as in t-tests), we compare timing distributions via their deciles. Let $F$ and $R$ denote timing samples from fixed (e.g., all-zeros) and random input classes, respectively. The test statistic is the vector of quantile differences:
\begin{equation}
\Delta = \hat{q}(F) - \hat{q}(R) \in \mathbb{R}^9
\end{equation}
where $\hat{q}_p$ denotes the empirical $p$-th quantile for $p \in \{0.1, 0.2, \ldots, 0.9\}$. This captures both uniform shifts (all quantiles affected equally, indicating branch changes) and tail effects (upper quantiles affected more, indicating cache-based leaks).

The decision functional is $m(\delta) = \max_k |\delta_k|$---the maximum absolute quantile difference. A timing leak exceeding threshold $\theta$ is declared when $m(\delta) > \theta$.

\paragraph{Likelihood.}
Quantile estimators are asymptotically normal, but real timing data can deviate from this approximation when covariance is underestimated. To prevent pathological certainty under covariance misspecification, we use a robust Student-$t$ likelihood via a scale-mixture representation:
\begin{align}
\Delta \mid \delta, \kappa &\sim \mathcal{N}\!\left(\delta, \frac{\Sigma_n}{\kappa}\right) \\
\kappa &\sim \text{Gamma}\!\left(\frac{\nu_\ell}{2}, \frac{\nu_\ell}{2}\right)
\end{align}
with $\nu_\ell = 8$. Marginally, this yields $\Delta \mid \delta \sim t_{\nu_\ell}(\delta, \Sigma_n)$.

When the estimated covariance $\Sigma_n$ is too small (e.g., due to unmodeled autocorrelation), $\kappa$ is pulled below 1, automatically inflating uncertainty. This is more robust than a Gaussian likelihood, which would produce overconfident posteriors under covariance misestimation.

\paragraph{Prior.}
We place a multivariate Student-$t$ prior on the latent effect profile $\delta$:
\begin{equation}
\delta \sim t_\nu(0, \sigma^2 R), \quad \nu = 4
\end{equation}
where $R = \text{Corr}(\Sigma_{\text{rate}})$ is the prior correlation matrix derived from the covariance rate estimated during calibration. The scale $\sigma$ is calibrated so that the prior probability of exceeding the effective threshold $\theta_{\text{eff}}$ equals a target $\pi_0 = 0.62$:
\begin{equation}
P\!\left(\max_k |\delta_k| > \theta_{\text{eff}} \mid \delta \sim t_\nu(0, \sigma^2 R)\right) = \pi_0
\end{equation}
This is solved by deterministic 1D root-finding using Monte Carlo integration (50,000 samples). The resulting prior expresses genuine uncertainty: it does not strongly favor either ``leak'' or ``no leak'' before seeing data.

We implement inference using the equivalent scale-mixture representation:
\begin{align}
\lambda &\sim \text{Gamma}\!\left(\frac{\nu}{2}, \frac{\nu}{2}\right) \\
\delta \mid \lambda &\sim \mathcal{N}\!\left(0, \frac{\sigma^2}{\lambda} R\right)
\end{align}
which enables closed-form Gibbs conditionals.

\subsection{Inference Procedure}
\label{app:inference}

\paragraph{Gibbs Sampler.}
We approximate the posterior $p(\delta \mid \Delta)$ using a deterministic Gibbs sampler over $(\delta, \lambda, \kappa)$. Each iteration samples:
\begin{enumerate}
\item $\delta^{(t)} \sim p(\delta \mid \lambda^{(t-1)}, \kappa^{(t-1)}, \Delta)$ --- a multivariate normal
\item $\lambda^{(t)} \sim p(\lambda \mid \delta^{(t)})$ --- a gamma distribution
\item $\kappa^{(t)} \sim p(\kappa \mid \delta^{(t)}, \Delta)$ --- a gamma distribution
\end{enumerate}

All conditionals are conjugate and admit closed-form sampling. The $\delta$ conditional requires computing a posterior precision matrix $Q = \kappa \Sigma_n^{-1} + (\lambda/\sigma^2) R^{-1}$ and sampling from $\mathcal{N}(Q^{-1}(\kappa \Sigma_n^{-1} \Delta), Q^{-1})$. All matrix operations use Cholesky decomposition for numerical stability---we never form explicit inverses.

\paragraph{Convergence and Cost.}
We run 256 total iterations with 64 burn-in, retaining 192 samples. This is sufficient because: (1) the model is low-dimensional ($d=9$) with conjugate structure, ensuring fast mixing; (2) initialization at $\lambda^{(0)} = \kappa^{(0)} = 1$ (the prior means) provides a good starting point; and (3) the decision (leak probability) depends only on whether $m(\delta^{(s)}) > \theta_{\text{eff}}$, which is robust to modest posterior approximation error.

We track $\lambda$ mixing via its coefficient of variation (CV) and effective sample size (ESS). If $\lambda$ CV $< 0.1$ or $\lambda$ ESS $< 20$, we flag a potential convergence issue in diagnostics, though this rarely occurs in practice.

Computational cost per posterior update is $O(d^3 + d^2 \cdot N_{\text{gibbs}})$, dominated by the Cholesky decompositions. With $d=9$ and $N_{\text{gibbs}}=256$, this takes $<1$ms on modern hardware---negligible compared to measurement time.

\paragraph{Leak Probability.}
The leak probability is estimated as the Monte Carlo fraction of posterior samples exceeding the threshold:
\begin{equation}
\widehat{P}(\text{leak} > \theta_{\text{eff}} \mid \Delta) = \frac{1}{N_{\text{keep}}} \sum_{s=1}^{N_{\text{keep}}} \mathbf{1}\!\left[m(\delta^{(s)}) > \theta_{\text{eff}}\right]
\end{equation}

\subsection{Block Bootstrap for Autocorrelation}
\label{app:bootstrap}

\paragraph{Why Standard Bootstrap Fails.}
Timing measurements exhibit autocorrelation: nearby samples are more similar than distant ones due to cache state, CPU frequency scaling, and OS scheduling. Standard bootstrap assumes i.i.d.\ samples, systematically underestimating variance and producing overconfident posteriors.

For example, if consecutive measurements have autocorrelation $\rho = 0.3$, the effective sample size may be 30--50\% lower than the nominal count. Ignoring this inflates the false positive rate---the main failure mode we aim to prevent.

\paragraph{Stream-Based Block Bootstrap.}
We use block bootstrap on the \emph{acquisition stream}---the interleaved sequence of (class, timing) pairs in measurement order---rather than per-class samples. This preserves the true temporal dependence structure.

The algorithm resamples contiguous blocks of length $\hat{b}$ with replacement, reconstructs a bootstrap stream of the same length, splits by class, and computes quantile differences. Over $B = 2000$ iterations, we estimate the covariance matrix $\Sigma_{\text{cal}}$ of $\Delta$.

\paragraph{Block Length Selection.}
We use the Politis-White automatic block length selection algorithm with a class-conditional modification. Standard Politis-White estimates autocorrelation from raw lag-$k$ pairs, but with interleaved sampling, a lag-$k$ pair may cross class boundaries. Instead, we compute autocorrelation using only same-class pairs at each acquisition-stream lag, then take the conservative maximum across classes.

The optimal block length scales as $\hat{b} \propto T^{1/3}$ where $T$ is the stream length. We clamp $\hat{b} \in [10, \min(3\sqrt{T}, T/3)]$ to prevent degenerate blocks. In high-autocorrelation regimes (detected when $\rho_{\text{max}}(10) > 0.3$), we inflate the block length by 50\% for additional conservatism.

\paragraph{Effective Sample Size via IACT.}
We compute the integrated autocorrelation time (IACT) $\hat{\tau}$ using the Geyer initial monotone sequence (IMS) algorithm. This estimates how many nominal samples equal one independent sample:
\begin{equation}
n_{\text{eff}} = \frac{n}{\hat{\tau}}
\end{equation}

For timing analysis, IACT is computed on indicator series $\mathbf{1}[y \leq q_p]$ for each quantile $p$, taking the maximum across quantiles and classes. This captures dependence in the quantile estimators, not just raw timings.

\paragraph{Covariance Scaling.}
The key insight enabling adaptive sampling is that quantile estimator covariance scales as $1/n$. We estimate a ``covariance rate'' during calibration:
\begin{equation}
\Sigma_{\text{rate}} = \hat{\Sigma}_{\text{cal}} \cdot n_{\text{cal}}
\end{equation}
During the adaptive loop with $n$ samples and current IACT $\hat{\tau}$:
\begin{equation}
\Sigma_n = \Sigma_{\text{rate}} \cdot \frac{\hat{\tau}}{\hat{\tau}_{\text{cal}}} \cdot \frac{1}{n}
\end{equation}
This allows cheap posterior updates without re-bootstrapping at each batch.

\subsection{Quality Gates}
\label{app:gates}

Quality gates detect when data quality is insufficient for a confident verdict, forcing an \textsc{Inconclusive} outcome rather than an unreliable Pass or Fail. This is the key mechanism ensuring calibrated verdicts.

\paragraph{Gate 1: Insufficient Information Gain.}
We compute the KL divergence between Gaussian approximations of the prior and posterior:
\begin{multline}
\text{KL}(p_{\text{post}} \| p_{\text{prior}}) = \frac{1}{2}\Big(
\text{tr}(\Lambda_0^{-1}\Lambda_{\text{post}}) \\
+ \mu_{\text{post}}^\top \Lambda_0^{-1} \mu_{\text{post}}
- d + \ln\frac{|\Lambda_0|}{|\Lambda_{\text{post}}|}
\Big)
\end{multline}
where $\Lambda_0 = 2\sigma^2 R$ is the marginal prior covariance for $\nu=4$, and $(\mu_{\text{post}}, \Lambda_{\text{post}})$ are estimated from Gibbs samples.

\textbf{DataTooNoisy}: If KL $< 0.7$ nats after calibration, the posterior has barely moved from the prior---the data is too noisy to distinguish signal from noise.

\textbf{NotLearning}: If the sum of inter-batch KL divergences falls below 0.001 for 5+ consecutive batches, the posterior has stopped updating. Additional samples will not help.

\paragraph{Gate 2: Would Take Too Long.}
We extrapolate time to decision based on current convergence rate. If projected time exceeds 10$\times$ the configured budget, we stop early with \textsc{Inconclusive} rather than waste resources on a test that cannot conclude.

\paragraph{Gate 3: Resource Budget Exceeded.}
If elapsed time exceeds the configured time budget, or total samples exceed the configured maximum, we stop with \textsc{Inconclusive} and report the current posterior probability.

\paragraph{Gate 4: Condition Drift Detected.}
The covariance rate $\Sigma_{\text{rate}}$ is estimated during calibration. If measurement conditions change during the adaptive loop, this estimate becomes invalid. We detect drift by comparing:
\begin{itemize}
\item Variance ratio: $\sigma_{\text{post}}^2 / \sigma_{\text{cal}}^2$ outside $[0.5, 2.0]$
\item Autocorrelation change: $|\rho_{\text{post}}(1) - \rho_{\text{cal}}(1)| > 0.3$
\item Mean drift: $|\mu_{\text{post}} - \mu_{\text{cal}}| / \sigma_{\text{cal}} > 3.0$
\end{itemize}
Any trigger yields \textsc{Inconclusive} with reason \textsc{ConditionsChanged}.

\paragraph{Likelihood Inflation Warning.}
When the posterior mean of $\kappa$ falls below 0.3, the likelihood covariance was effectively scaled up by $>3\times$ for robustness. This indicates potential covariance misestimation and is reported in diagnostics, though it does not block verdicts.

\subsection{Measurement Floor Derivation}
\label{app:floor}

A critical design element is distinguishing what the user \emph{wants} to detect ($\theta_{\text{user}}$) from what the measurement \emph{can} detect ($\theta_{\text{floor}}$).

\paragraph{Floor-Rate Constant.}
During calibration, we compute a floor-rate constant $c_{\text{floor}}$ representing the 95th percentile of the null distribution of $m(\delta)$:
\begin{equation}
c_{\text{floor}} = q_{0.95}\!\left(\max_k |Z_{0,k}|\right), \quad Z_0 \sim \mathcal{N}(0, \Sigma_{\text{rate}})
\end{equation}
estimated via Monte Carlo (50,000 samples). This characterizes the ``noise floor'' of the measurement system.

\paragraph{Dynamic Measurement Floor.}
During the adaptive loop, the statistical floor shrinks with effective sample size:
\begin{equation}
\theta_{\text{floor,stat}}(n) = \frac{c_{\text{floor}}}{\sqrt{n_{\text{eff}}}}
\end{equation}

The timer also imposes a floor based on quantization:
\begin{equation}
\theta_{\text{tick}} = \frac{\text{1 tick (ns)}}{K}
\end{equation}
where $K$ is the batch size (if batching is used). The combined floor is:
\begin{equation}
\theta_{\text{floor}}(n) = \max\!\left(\theta_{\text{floor,stat}}(n), \theta_{\text{tick}}\right)
\end{equation}

\paragraph{Effective Threshold.}
The effective threshold used for inference is:
\begin{equation}
\theta_{\text{eff}} = \max(\theta_{\text{user}}, \theta_{\text{floor}})
\end{equation}

When $\theta_{\text{eff}} > \theta_{\text{user}}$:
\begin{itemize}
\item \textbf{Fail propagates}: Detecting $m(\delta) > \theta_{\text{eff}}$ implies $m(\delta) > \theta_{\text{user}}$, since $\theta_{\text{eff}} \geq \theta_{\text{user}}$.
\item \textbf{Pass does not propagate}: ``No effect above $\theta_{\text{eff}}$'' is compatible with effects in $(\theta_{\text{user}}, \theta_{\text{eff}}]$. We return \textsc{Inconclusive} with reason \textsc{ThresholdElevated}.
\end{itemize}

As sampling continues, $\theta_{\text{floor}}(n)$ decreases. If it drops to $\theta_{\text{user}}$ or below, Pass becomes possible (subject to the posterior). This is why adaptive sampling can eventually reach Pass even when early calibration shows an elevated floor.

\vspace{1em}
\noindent\textit{For complete algorithmic details including pseudocode, numerical stability procedures, and platform-specific considerations, see the specification document in the artifact repository.}
 in the main document

\section{Statistical Methodology Details}
\label{app:methodology}

This appendix provides additional detail on tacet's statistical methodology for reviewers familiar with Bayesian inference. We cover the probabilistic model, inference procedure, autocorrelation handling, quality gates, and measurement floor derivation. For complete algorithmic specifications, see the specification document in the artifact repository.

\subsection{Probabilistic Model}
\label{app:model}

\paragraph{Test Statistic.}
Rather than comparing means (as in t-tests), we compare timing distributions via their deciles. Let $F$ and $R$ denote timing samples from fixed (e.g., all-zeros) and random input classes, respectively. The test statistic is the vector of quantile differences:
\begin{equation}
\Delta = \hat{q}(F) - \hat{q}(R) \in \mathbb{R}^9
\end{equation}
where $\hat{q}_p$ denotes the empirical $p$-th quantile for $p \in \{0.1, 0.2, \ldots, 0.9\}$. This captures both uniform shifts (all quantiles affected equally, indicating branch changes) and tail effects (upper quantiles affected more, indicating cache-based leaks).

The decision functional is $m(\delta) = \max_k |\delta_k|$---the maximum absolute quantile difference. A timing leak exceeding threshold $\theta$ is declared when $m(\delta) > \theta$.

\paragraph{Likelihood.}
Quantile estimators are asymptotically normal, but real timing data can deviate from this approximation when covariance is underestimated. To prevent pathological certainty under covariance misspecification, we use a robust Student-$t$ likelihood via a scale-mixture representation:
\begin{align}
\Delta \mid \delta, \kappa &\sim \mathcal{N}\!\left(\delta, \frac{\Sigma_n}{\kappa}\right) \\
\kappa &\sim \text{Gamma}\!\left(\frac{\nu_\ell}{2}, \frac{\nu_\ell}{2}\right)
\end{align}
with $\nu_\ell = 8$. Marginally, this yields $\Delta \mid \delta \sim t_{\nu_\ell}(\delta, \Sigma_n)$.

When the estimated covariance $\Sigma_n$ is too small (e.g., due to unmodeled autocorrelation), $\kappa$ is pulled below 1, automatically inflating uncertainty. This is more robust than a Gaussian likelihood, which would produce overconfident posteriors under covariance misestimation.

\paragraph{Prior.}
We place a multivariate Student-$t$ prior on the latent effect profile $\delta$:
\begin{equation}
\delta \sim t_\nu(0, \sigma^2 R), \quad \nu = 4
\end{equation}
where $R = \text{Corr}(\Sigma_{\text{rate}})$ is the prior correlation matrix derived from the covariance rate estimated during calibration. The scale $\sigma$ is calibrated so that the prior probability of exceeding the effective threshold $\theta_{\text{eff}}$ equals a target $\pi_0 = 0.62$:
\begin{equation}
P\!\left(\max_k |\delta_k| > \theta_{\text{eff}} \mid \delta \sim t_\nu(0, \sigma^2 R)\right) = \pi_0
\end{equation}
This is solved by deterministic 1D root-finding using Monte Carlo integration (50,000 samples). The resulting prior expresses genuine uncertainty: it does not strongly favor either ``leak'' or ``no leak'' before seeing data.

We implement inference using the equivalent scale-mixture representation:
\begin{align}
\lambda &\sim \text{Gamma}\!\left(\frac{\nu}{2}, \frac{\nu}{2}\right) \\
\delta \mid \lambda &\sim \mathcal{N}\!\left(0, \frac{\sigma^2}{\lambda} R\right)
\end{align}
which enables closed-form Gibbs conditionals.

\subsection{Inference Procedure}
\label{app:inference}

\paragraph{Gibbs Sampler.}
We approximate the posterior $p(\delta \mid \Delta)$ using a deterministic Gibbs sampler over $(\delta, \lambda, \kappa)$. Each iteration samples:
\begin{enumerate}
\item $\delta^{(t)} \sim p(\delta \mid \lambda^{(t-1)}, \kappa^{(t-1)}, \Delta)$ --- a multivariate normal
\item $\lambda^{(t)} \sim p(\lambda \mid \delta^{(t)})$ --- a gamma distribution
\item $\kappa^{(t)} \sim p(\kappa \mid \delta^{(t)}, \Delta)$ --- a gamma distribution
\end{enumerate}

All conditionals are conjugate and admit closed-form sampling. The $\delta$ conditional requires computing a posterior precision matrix $Q = \kappa \Sigma_n^{-1} + (\lambda/\sigma^2) R^{-1}$ and sampling from $\mathcal{N}(Q^{-1}(\kappa \Sigma_n^{-1} \Delta), Q^{-1})$. All matrix operations use Cholesky decomposition for numerical stability---we never form explicit inverses.

\paragraph{Convergence and Cost.}
We run 256 total iterations with 64 burn-in, retaining 192 samples. This is sufficient because: (1) the model is low-dimensional ($d=9$) with conjugate structure, ensuring fast mixing; (2) initialization at $\lambda^{(0)} = \kappa^{(0)} = 1$ (the prior means) provides a good starting point; and (3) the decision (leak probability) depends only on whether $m(\delta^{(s)}) > \theta_{\text{eff}}$, which is robust to modest posterior approximation error.

We track $\lambda$ mixing via its coefficient of variation (CV) and effective sample size (ESS). If $\lambda$ CV $< 0.1$ or $\lambda$ ESS $< 20$, we flag a potential convergence issue in diagnostics, though this rarely occurs in practice.

Computational cost per posterior update is $O(d^3 + d^2 \cdot N_{\text{gibbs}})$, dominated by the Cholesky decompositions. With $d=9$ and $N_{\text{gibbs}}=256$, this takes $<1$ms on modern hardware---negligible compared to measurement time.

\paragraph{Leak Probability.}
The leak probability is estimated as the Monte Carlo fraction of posterior samples exceeding the threshold:
\begin{equation}
\widehat{P}(\text{leak} > \theta_{\text{eff}} \mid \Delta) = \frac{1}{N_{\text{keep}}} \sum_{s=1}^{N_{\text{keep}}} \mathbf{1}\!\left[m(\delta^{(s)}) > \theta_{\text{eff}}\right]
\end{equation}

\subsection{Block Bootstrap for Autocorrelation}
\label{app:bootstrap}

\paragraph{Why Standard Bootstrap Fails.}
Timing measurements exhibit autocorrelation: nearby samples are more similar than distant ones due to cache state, CPU frequency scaling, and OS scheduling. Standard bootstrap assumes i.i.d.\ samples, systematically underestimating variance and producing overconfident posteriors.

For example, if consecutive measurements have autocorrelation $\rho = 0.3$, the effective sample size may be 30--50\% lower than the nominal count. Ignoring this inflates the false positive rate---the main failure mode we aim to prevent.

\paragraph{Stream-Based Block Bootstrap.}
We use block bootstrap on the \emph{acquisition stream}---the interleaved sequence of (class, timing) pairs in measurement order---rather than per-class samples. This preserves the true temporal dependence structure.

The algorithm resamples contiguous blocks of length $\hat{b}$ with replacement, reconstructs a bootstrap stream of the same length, splits by class, and computes quantile differences. Over $B = 2000$ iterations, we estimate the covariance matrix $\Sigma_{\text{cal}}$ of $\Delta$.

\paragraph{Block Length Selection.}
We use the Politis-White automatic block length selection algorithm with a class-conditional modification. Standard Politis-White estimates autocorrelation from raw lag-$k$ pairs, but with interleaved sampling, a lag-$k$ pair may cross class boundaries. Instead, we compute autocorrelation using only same-class pairs at each acquisition-stream lag, then take the conservative maximum across classes.

The optimal block length scales as $\hat{b} \propto T^{1/3}$ where $T$ is the stream length. We clamp $\hat{b} \in [10, \min(3\sqrt{T}, T/3)]$ to prevent degenerate blocks. In high-autocorrelation regimes (detected when $\rho_{\text{max}}(10) > 0.3$), we inflate the block length by 50\% for additional conservatism.

\paragraph{Effective Sample Size via IACT.}
We compute the integrated autocorrelation time (IACT) $\hat{\tau}$ using the Geyer initial monotone sequence (IMS) algorithm. This estimates how many nominal samples equal one independent sample:
\begin{equation}
n_{\text{eff}} = \frac{n}{\hat{\tau}}
\end{equation}

For timing analysis, IACT is computed on indicator series $\mathbf{1}[y \leq q_p]$ for each quantile $p$, taking the maximum across quantiles and classes. This captures dependence in the quantile estimators, not just raw timings.

\paragraph{Covariance Scaling.}
The key insight enabling adaptive sampling is that quantile estimator covariance scales as $1/n$. We estimate a ``covariance rate'' during calibration:
\begin{equation}
\Sigma_{\text{rate}} = \hat{\Sigma}_{\text{cal}} \cdot n_{\text{cal}}
\end{equation}
During the adaptive loop with $n$ samples and current IACT $\hat{\tau}$:
\begin{equation}
\Sigma_n = \Sigma_{\text{rate}} \cdot \frac{\hat{\tau}}{\hat{\tau}_{\text{cal}}} \cdot \frac{1}{n}
\end{equation}
This allows cheap posterior updates without re-bootstrapping at each batch.

\subsection{Quality Gates}
\label{app:gates}

Quality gates detect when data quality is insufficient for a confident verdict, forcing an \textsc{Inconclusive} outcome rather than an unreliable Pass or Fail. This is the key mechanism ensuring calibrated verdicts.

\paragraph{Gate 1: Insufficient Information Gain.}
We compute the KL divergence between Gaussian approximations of the prior and posterior:
\begin{multline}
\text{KL}(p_{\text{post}} \| p_{\text{prior}}) = \frac{1}{2}\Big(
\text{tr}(\Lambda_0^{-1}\Lambda_{\text{post}}) \\
+ \mu_{\text{post}}^\top \Lambda_0^{-1} \mu_{\text{post}}
- d + \ln\frac{|\Lambda_0|}{|\Lambda_{\text{post}}|}
\Big)
\end{multline}
where $\Lambda_0 = 2\sigma^2 R$ is the marginal prior covariance for $\nu=4$, and $(\mu_{\text{post}}, \Lambda_{\text{post}})$ are estimated from Gibbs samples.

\textbf{DataTooNoisy}: If KL $< 0.7$ nats after calibration, the posterior has barely moved from the prior---the data is too noisy to distinguish signal from noise.

\textbf{NotLearning}: If the sum of inter-batch KL divergences falls below 0.001 for 5+ consecutive batches, the posterior has stopped updating. Additional samples will not help.

\paragraph{Gate 2: Would Take Too Long.}
We extrapolate time to decision based on current convergence rate. If projected time exceeds 10$\times$ the configured budget, we stop early with \textsc{Inconclusive} rather than waste resources on a test that cannot conclude.

\paragraph{Gate 3: Resource Budget Exceeded.}
If elapsed time exceeds the configured time budget, or total samples exceed the configured maximum, we stop with \textsc{Inconclusive} and report the current posterior probability.

\paragraph{Gate 4: Condition Drift Detected.}
The covariance rate $\Sigma_{\text{rate}}$ is estimated during calibration. If measurement conditions change during the adaptive loop, this estimate becomes invalid. We detect drift by comparing:
\begin{itemize}
\item Variance ratio: $\sigma_{\text{post}}^2 / \sigma_{\text{cal}}^2$ outside $[0.5, 2.0]$
\item Autocorrelation change: $|\rho_{\text{post}}(1) - \rho_{\text{cal}}(1)| > 0.3$
\item Mean drift: $|\mu_{\text{post}} - \mu_{\text{cal}}| / \sigma_{\text{cal}} > 3.0$
\end{itemize}
Any trigger yields \textsc{Inconclusive} with reason \textsc{ConditionsChanged}.

\paragraph{Likelihood Inflation Warning.}
When the posterior mean of $\kappa$ falls below 0.3, the likelihood covariance was effectively scaled up by $>3\times$ for robustness. This indicates potential covariance misestimation and is reported in diagnostics, though it does not block verdicts.

\subsection{Measurement Floor Derivation}
\label{app:floor}

A critical design element is distinguishing what the user \emph{wants} to detect ($\theta_{\text{user}}$) from what the measurement \emph{can} detect ($\theta_{\text{floor}}$).

\paragraph{Floor-Rate Constant.}
During calibration, we compute a floor-rate constant $c_{\text{floor}}$ representing the 95th percentile of the null distribution of $m(\delta)$:
\begin{equation}
c_{\text{floor}} = q_{0.95}\!\left(\max_k |Z_{0,k}|\right), \quad Z_0 \sim \mathcal{N}(0, \Sigma_{\text{rate}})
\end{equation}
estimated via Monte Carlo (50,000 samples). This characterizes the ``noise floor'' of the measurement system.

\paragraph{Dynamic Measurement Floor.}
During the adaptive loop, the statistical floor shrinks with effective sample size:
\begin{equation}
\theta_{\text{floor,stat}}(n) = \frac{c_{\text{floor}}}{\sqrt{n_{\text{eff}}}}
\end{equation}

The timer also imposes a floor based on quantization:
\begin{equation}
\theta_{\text{tick}} = \frac{\text{1 tick (ns)}}{K}
\end{equation}
where $K$ is the batch size (if batching is used). The combined floor is:
\begin{equation}
\theta_{\text{floor}}(n) = \max\!\left(\theta_{\text{floor,stat}}(n), \theta_{\text{tick}}\right)
\end{equation}

\paragraph{Effective Threshold.}
The effective threshold used for inference is:
\begin{equation}
\theta_{\text{eff}} = \max(\theta_{\text{user}}, \theta_{\text{floor}})
\end{equation}

When $\theta_{\text{eff}} > \theta_{\text{user}}$:
\begin{itemize}
\item \textbf{Fail propagates}: Detecting $m(\delta) > \theta_{\text{eff}}$ implies $m(\delta) > \theta_{\text{user}}$, since $\theta_{\text{eff}} \geq \theta_{\text{user}}$.
\item \textbf{Pass does not propagate}: ``No effect above $\theta_{\text{eff}}$'' is compatible with effects in $(\theta_{\text{user}}, \theta_{\text{eff}}]$. We return \textsc{Inconclusive} with reason \textsc{ThresholdElevated}.
\end{itemize}

As sampling continues, $\theta_{\text{floor}}(n)$ decreases. If it drops to $\theta_{\text{user}}$ or below, Pass becomes possible (subject to the posterior). This is why adaptive sampling can eventually reach Pass even when early calibration shows an elevated floor.

\vspace{1em}
\noindent\textit{For complete algorithmic details including pseudocode, numerical stability procedures, and platform-specific considerations, see the specification document in the artifact repository.}
 in the main document

\section{Statistical Methodology Details}
\label{app:methodology}

This appendix provides additional detail on tacet's statistical methodology for reviewers familiar with Bayesian inference. We cover the probabilistic model, inference procedure, autocorrelation handling, quality gates, and measurement floor derivation. For complete algorithmic specifications, see the specification document in the artifact repository.

\subsection{Probabilistic Model}
\label{app:model}

\paragraph{Test Statistic.}
Rather than comparing means (as in t-tests), we compare timing distributions via their deciles. Let $F$ and $R$ denote timing samples from fixed (e.g., all-zeros) and random input classes, respectively. The test statistic is the vector of quantile differences:
\begin{equation}
\Delta = \hat{q}(F) - \hat{q}(R) \in \mathbb{R}^9
\end{equation}
where $\hat{q}_p$ denotes the empirical $p$-th quantile for $p \in \{0.1, 0.2, \ldots, 0.9\}$. This captures both uniform shifts (all quantiles affected equally, indicating branch changes) and tail effects (upper quantiles affected more, indicating cache-based leaks).

The decision functional is $m(\delta) = \max_k |\delta_k|$---the maximum absolute quantile difference. A timing leak exceeding threshold $\theta$ is declared when $m(\delta) > \theta$.

\paragraph{Likelihood.}
Quantile estimators are asymptotically normal, but real timing data can deviate from this approximation when covariance is underestimated. To prevent pathological certainty under covariance misspecification, we use a robust Student-$t$ likelihood via a scale-mixture representation:
\begin{align}
\Delta \mid \delta, \kappa &\sim \mathcal{N}\!\left(\delta, \frac{\Sigma_n}{\kappa}\right) \\
\kappa &\sim \text{Gamma}\!\left(\frac{\nu_\ell}{2}, \frac{\nu_\ell}{2}\right)
\end{align}
with $\nu_\ell = 8$. Marginally, this yields $\Delta \mid \delta \sim t_{\nu_\ell}(\delta, \Sigma_n)$.

When the estimated covariance $\Sigma_n$ is too small (e.g., due to unmodeled autocorrelation), $\kappa$ is pulled below 1, automatically inflating uncertainty. This is more robust than a Gaussian likelihood, which would produce overconfident posteriors under covariance misestimation.

\paragraph{Prior.}
We place a multivariate Student-$t$ prior on the latent effect profile $\delta$:
\begin{equation}
\delta \sim t_\nu(0, \sigma^2 R), \quad \nu = 4
\end{equation}
where $R = \text{Corr}(\Sigma_{\text{rate}})$ is the prior correlation matrix derived from the covariance rate estimated during calibration. The scale $\sigma$ is calibrated so that the prior probability of exceeding the effective threshold $\theta_{\text{eff}}$ equals a target $\pi_0 = 0.62$:
\begin{equation}
P\!\left(\max_k |\delta_k| > \theta_{\text{eff}} \mid \delta \sim t_\nu(0, \sigma^2 R)\right) = \pi_0
\end{equation}
This is solved by deterministic 1D root-finding using Monte Carlo integration (50,000 samples). The resulting prior expresses genuine uncertainty: it does not strongly favor either ``leak'' or ``no leak'' before seeing data.

We implement inference using the equivalent scale-mixture representation:
\begin{align}
\lambda &\sim \text{Gamma}\!\left(\frac{\nu}{2}, \frac{\nu}{2}\right) \\
\delta \mid \lambda &\sim \mathcal{N}\!\left(0, \frac{\sigma^2}{\lambda} R\right)
\end{align}
which enables closed-form Gibbs conditionals.

\subsection{Inference Procedure}
\label{app:inference}

\paragraph{Gibbs Sampler.}
We approximate the posterior $p(\delta \mid \Delta)$ using a deterministic Gibbs sampler over $(\delta, \lambda, \kappa)$. Each iteration samples:
\begin{enumerate}
\item $\delta^{(t)} \sim p(\delta \mid \lambda^{(t-1)}, \kappa^{(t-1)}, \Delta)$ --- a multivariate normal
\item $\lambda^{(t)} \sim p(\lambda \mid \delta^{(t)})$ --- a gamma distribution
\item $\kappa^{(t)} \sim p(\kappa \mid \delta^{(t)}, \Delta)$ --- a gamma distribution
\end{enumerate}

All conditionals are conjugate and admit closed-form sampling. The $\delta$ conditional requires computing a posterior precision matrix $Q = \kappa \Sigma_n^{-1} + (\lambda/\sigma^2) R^{-1}$ and sampling from $\mathcal{N}(Q^{-1}(\kappa \Sigma_n^{-1} \Delta), Q^{-1})$. All matrix operations use Cholesky decomposition for numerical stability---we never form explicit inverses.

\paragraph{Convergence and Cost.}
We run 256 total iterations with 64 burn-in, retaining 192 samples. This is sufficient because: (1) the model is low-dimensional ($d=9$) with conjugate structure, ensuring fast mixing; (2) initialization at $\lambda^{(0)} = \kappa^{(0)} = 1$ (the prior means) provides a good starting point; and (3) the decision (leak probability) depends only on whether $m(\delta^{(s)}) > \theta_{\text{eff}}$, which is robust to modest posterior approximation error.

We track $\lambda$ mixing via its coefficient of variation (CV) and effective sample size (ESS). If $\lambda$ CV $< 0.1$ or $\lambda$ ESS $< 20$, we flag a potential convergence issue in diagnostics, though this rarely occurs in practice.

Computational cost per posterior update is $O(d^3 + d^2 \cdot N_{\text{gibbs}})$, dominated by the Cholesky decompositions. With $d=9$ and $N_{\text{gibbs}}=256$, this takes $<1$ms on modern hardware---negligible compared to measurement time.

\paragraph{Leak Probability.}
The leak probability is estimated as the Monte Carlo fraction of posterior samples exceeding the threshold:
\begin{equation}
\widehat{P}(\text{leak} > \theta_{\text{eff}} \mid \Delta) = \frac{1}{N_{\text{keep}}} \sum_{s=1}^{N_{\text{keep}}} \mathbf{1}\!\left[m(\delta^{(s)}) > \theta_{\text{eff}}\right]
\end{equation}

\subsection{Block Bootstrap for Autocorrelation}
\label{app:bootstrap}

\paragraph{Why Standard Bootstrap Fails.}
Timing measurements exhibit autocorrelation: nearby samples are more similar than distant ones due to cache state, CPU frequency scaling, and OS scheduling. Standard bootstrap assumes i.i.d.\ samples, systematically underestimating variance and producing overconfident posteriors.

For example, if consecutive measurements have autocorrelation $\rho = 0.3$, the effective sample size may be 30--50\% lower than the nominal count. Ignoring this inflates the false positive rate---the main failure mode we aim to prevent.

\paragraph{Stream-Based Block Bootstrap.}
We use block bootstrap on the \emph{acquisition stream}---the interleaved sequence of (class, timing) pairs in measurement order---rather than per-class samples. This preserves the true temporal dependence structure.

The algorithm resamples contiguous blocks of length $\hat{b}$ with replacement, reconstructs a bootstrap stream of the same length, splits by class, and computes quantile differences. Over $B = 2000$ iterations, we estimate the covariance matrix $\Sigma_{\text{cal}}$ of $\Delta$.

\paragraph{Block Length Selection.}
We use the Politis-White automatic block length selection algorithm with a class-conditional modification. Standard Politis-White estimates autocorrelation from raw lag-$k$ pairs, but with interleaved sampling, a lag-$k$ pair may cross class boundaries. Instead, we compute autocorrelation using only same-class pairs at each acquisition-stream lag, then take the conservative maximum across classes.

The optimal block length scales as $\hat{b} \propto T^{1/3}$ where $T$ is the stream length. We clamp $\hat{b} \in [10, \min(3\sqrt{T}, T/3)]$ to prevent degenerate blocks. In high-autocorrelation regimes (detected when $\rho_{\text{max}}(10) > 0.3$), we inflate the block length by 50\% for additional conservatism.

\paragraph{Effective Sample Size via IACT.}
We compute the integrated autocorrelation time (IACT) $\hat{\tau}$ using the Geyer initial monotone sequence (IMS) algorithm. This estimates how many nominal samples equal one independent sample:
\begin{equation}
n_{\text{eff}} = \frac{n}{\hat{\tau}}
\end{equation}

For timing analysis, IACT is computed on indicator series $\mathbf{1}[y \leq q_p]$ for each quantile $p$, taking the maximum across quantiles and classes. This captures dependence in the quantile estimators, not just raw timings.

\paragraph{Covariance Scaling.}
The key insight enabling adaptive sampling is that quantile estimator covariance scales as $1/n$. We estimate a ``covariance rate'' during calibration:
\begin{equation}
\Sigma_{\text{rate}} = \hat{\Sigma}_{\text{cal}} \cdot n_{\text{cal}}
\end{equation}
During the adaptive loop with $n$ samples and current IACT $\hat{\tau}$:
\begin{equation}
\Sigma_n = \Sigma_{\text{rate}} \cdot \frac{\hat{\tau}}{\hat{\tau}_{\text{cal}}} \cdot \frac{1}{n}
\end{equation}
This allows cheap posterior updates without re-bootstrapping at each batch.

\subsection{Quality Gates}
\label{app:gates}

Quality gates detect when data quality is insufficient for a confident verdict, forcing an \textsc{Inconclusive} outcome rather than an unreliable Pass or Fail. This is the key mechanism ensuring calibrated verdicts.

\paragraph{Gate 1: Insufficient Information Gain.}
We compute the KL divergence between Gaussian approximations of the prior and posterior:
\begin{multline}
\text{KL}(p_{\text{post}} \| p_{\text{prior}}) = \frac{1}{2}\Big(
\text{tr}(\Lambda_0^{-1}\Lambda_{\text{post}}) \\
+ \mu_{\text{post}}^\top \Lambda_0^{-1} \mu_{\text{post}}
- d + \ln\frac{|\Lambda_0|}{|\Lambda_{\text{post}}|}
\Big)
\end{multline}
where $\Lambda_0 = 2\sigma^2 R$ is the marginal prior covariance for $\nu=4$, and $(\mu_{\text{post}}, \Lambda_{\text{post}})$ are estimated from Gibbs samples.

\textbf{DataTooNoisy}: If KL $< 0.7$ nats after calibration, the posterior has barely moved from the prior---the data is too noisy to distinguish signal from noise.

\textbf{NotLearning}: If the sum of inter-batch KL divergences falls below 0.001 for 5+ consecutive batches, the posterior has stopped updating. Additional samples will not help.

\paragraph{Gate 2: Would Take Too Long.}
We extrapolate time to decision based on current convergence rate. If projected time exceeds 10$\times$ the configured budget, we stop early with \textsc{Inconclusive} rather than waste resources on a test that cannot conclude.

\paragraph{Gate 3: Resource Budget Exceeded.}
If elapsed time exceeds the configured time budget, or total samples exceed the configured maximum, we stop with \textsc{Inconclusive} and report the current posterior probability.

\paragraph{Gate 4: Condition Drift Detected.}
The covariance rate $\Sigma_{\text{rate}}$ is estimated during calibration. If measurement conditions change during the adaptive loop, this estimate becomes invalid. We detect drift by comparing:
\begin{itemize}
\item Variance ratio: $\sigma_{\text{post}}^2 / \sigma_{\text{cal}}^2$ outside $[0.5, 2.0]$
\item Autocorrelation change: $|\rho_{\text{post}}(1) - \rho_{\text{cal}}(1)| > 0.3$
\item Mean drift: $|\mu_{\text{post}} - \mu_{\text{cal}}| / \sigma_{\text{cal}} > 3.0$
\end{itemize}
Any trigger yields \textsc{Inconclusive} with reason \textsc{ConditionsChanged}.

\paragraph{Likelihood Inflation Warning.}
When the posterior mean of $\kappa$ falls below 0.3, the likelihood covariance was effectively scaled up by $>3\times$ for robustness. This indicates potential covariance misestimation and is reported in diagnostics, though it does not block verdicts.

\subsection{Measurement Floor Derivation}
\label{app:floor}

A critical design element is distinguishing what the user \emph{wants} to detect ($\theta_{\text{user}}$) from what the measurement \emph{can} detect ($\theta_{\text{floor}}$).

\paragraph{Floor-Rate Constant.}
During calibration, we compute a floor-rate constant $c_{\text{floor}}$ representing the 95th percentile of the null distribution of $m(\delta)$:
\begin{equation}
c_{\text{floor}} = q_{0.95}\!\left(\max_k |Z_{0,k}|\right), \quad Z_0 \sim \mathcal{N}(0, \Sigma_{\text{rate}})
\end{equation}
estimated via Monte Carlo (50,000 samples). This characterizes the ``noise floor'' of the measurement system.

\paragraph{Dynamic Measurement Floor.}
During the adaptive loop, the statistical floor shrinks with effective sample size:
\begin{equation}
\theta_{\text{floor,stat}}(n) = \frac{c_{\text{floor}}}{\sqrt{n_{\text{eff}}}}
\end{equation}

The timer also imposes a floor based on quantization:
\begin{equation}
\theta_{\text{tick}} = \frac{\text{1 tick (ns)}}{K}
\end{equation}
where $K$ is the batch size (if batching is used). The combined floor is:
\begin{equation}
\theta_{\text{floor}}(n) = \max\!\left(\theta_{\text{floor,stat}}(n), \theta_{\text{tick}}\right)
\end{equation}

\paragraph{Effective Threshold.}
The effective threshold used for inference is:
\begin{equation}
\theta_{\text{eff}} = \max(\theta_{\text{user}}, \theta_{\text{floor}})
\end{equation}

When $\theta_{\text{eff}} > \theta_{\text{user}}$:
\begin{itemize}
\item \textbf{Fail propagates}: Detecting $m(\delta) > \theta_{\text{eff}}$ implies $m(\delta) > \theta_{\text{user}}$, since $\theta_{\text{eff}} \geq \theta_{\text{user}}$.
\item \textbf{Pass does not propagate}: ``No effect above $\theta_{\text{eff}}$'' is compatible with effects in $(\theta_{\text{user}}, \theta_{\text{eff}}]$. We return \textsc{Inconclusive} with reason \textsc{ThresholdElevated}.
\end{itemize}

As sampling continues, $\theta_{\text{floor}}(n)$ decreases. If it drops to $\theta_{\text{user}}$ or below, Pass becomes possible (subject to the posterior). This is why adaptive sampling can eventually reach Pass even when early calibration shows an elevated floor.

\vspace{1em}
\noindent\textit{For complete algorithmic details including pseudocode, numerical stability procedures, and platform-specific considerations, see the specification document in the artifact repository.}
